\documentclass[11pt,a4paper]{article}
\usepackage[T1]{fontenc}
\usepackage{lmodern,amsmath,amssymb,amsthm,booktabs,microtype}
\usepackage[margin=26mm]{geometry}
\usepackage[colorlinks=true,linkcolor=blue,citecolor=blue,urlcolor=blue]{hyperref}
\newtheorem{theorem}{Theorem}[section]
\newtheorem{lemma}[theorem]{Lemma}
\newtheorem{proposition}[theorem]{Proposition}
\newcommand{\Q}{\mathbb Q}\newcommand{\Z}{\mathbb Z}
\newcommand{\Tr}{\operatorname{Tr}}
\setlength{\emergencystretch}{2em}
\title{Three-detector quartics with 229 rational coordinates}
\author{}\date{17 November 2026}
\begin{document}\maketitle

\begin{abstract}
A complete quadratic detector can be duplicated using 34 additional
rational coordinates. Both components of the elliptic abscissa sum
are already quadratic expressions in the retained coordinates.
Consequently three additive detector arguments require 229 total
coordinates and 207 quadratic equations. The resulting primitive
quartic has 3,312 monomials and exact coefficient height
$\max\{9063965991649194484576286065455625302,149^{2N+2}\}$.
Every fixed $N\ge15$ defines the integers under the stated arithmetic
premises. At $N=15$ its coefficient height is $149^{32}$, a 232-bit
integer. The argument shifts remain uniformly separated on the full
finite elliptic point domain, including its torsion strata.
\end{abstract}

\section{Exact specification and arithmetic scope}

Every count includes the distinguished input $X=x_1$.
Fix $N\ge1$ and the integer coefficient $C=149^N$.
Let $\mathcal S_{195}(C)$ be the 175-row system of~\cite{Parent195}.
Its complete quadratic table has SHA-256 digest
\begin{center}\small\ttfamily
304f94843c1a4ac94fcd5f1295fee8b532e7f0424e638314d892c556e2fbfdf9.
\end{center}
All source reference, scalar, comparison and elliptic equations are
retained. Their interpretations use the exact source arithmetic
premises and the certified source coordinate transformations.
These premises include the elliptic group description, scalar
contact-height argument, pole obstruction and prepared realization,
modular contact-height estimates, and fixed-data simultaneous-prime
theorem. Finite coefficient verification does not replace these
arithmetic results.

In the source table the second detector has precisely the following
34 local coordinates:
\[
 \{7,8,9,10\}\ \cup\ \{86,\ldots,99\}
 \ \cup\ \{134,\ldots,149\}.
\]
In increasing order, copy these coordinates to
$x_{196},\ldots,x_{229}$. Copy each of the 31 source rows whose
name contains \texttt{detector1} or \texttt{det1}, replacing just
these local coordinates by their copies. In the names replace the
detector index by two. The newly stored argument is $x_{200}$.
Replace the contact equations of all three detectors by the exact
center-$3/11$ circuit below; the shared point and scalar coordinates
are unchanged.

Only the argument $x_{86}$ has a consumer outside the source second
detector: the final aggregate row. Thus this operation is a copy of
the full detector, including its contact circuit, two trace
components, pole realization, auxiliary squares, and nonzero guard.
It does not omit an external recovery equation or require a new
shared coordinate. Replace the source aggregate row and add the
second component row specified in the next section. This gives
$175+31+1=207$ rows in $195+34=229$ coordinates.

The files \path{quadratic.json}, \path{polynomial.json}, and
\path{graft.json} give, respectively, all 207 integer quadrics,
their complete quartic expansion, and the coordinate and row maps.
The only fixed parameter in these tables is $C$, which is not
quantified. The directory
\path{computation/quartic/explicit229_v62/} contains these files and
\path{contact-change.json}, which specifies every changed contact row.

\begin{theorem}\label{thm:main}
Under the stated arithmetic inputs, for every integer $N\ge15$ the
specified quartic $P_N$ satisfies
\[
 X\in\Z\quad\Longleftrightarrow\quad
 \exists y_2,\ldots,y_{229}\in\Q:\quad
 P_N(X,y_2,\ldots,y_{229})=0.
\]
It is primitive, has degree four and 3,312 monomials, and at $N=15$
has exact coefficient height $149^{32}$. Every helper and guard is
included in the 229-coordinate count.
\end{theorem}

\section{Both abscissa components without a recovery coordinate}

Put $\rho=\sqrt2$, $d=3-\rho$, $D=d^2=11-6\rho$ and
$E:y^2=x^3-Dx$. Set
\[
 R_0=(1+\rho,2\rho),\qquad P_0=2R_0.
\]
Let $Q_1,Q_2,Q_3,Q_4$ be the finite elliptic
points recovered by the source reciprocal and inverse Vieta charts.
Write
\[
 \mathcal S=\sum_{i=1}^4 x(Q_i)=S_0+S_1\rho.
\]
The source permits finite torsion points in the first, third, and
fourth difference blocks. The scalar inverse Vieta point $Q_2$
is nontorsion. The following formulas do not require a difference
point to be nontorsion.

Let $q=x_{77}$ and $w=x_{78}$. Their literal guard is
$qw=1$. In the final source coordinates set
\begin{align*}
 f_1&=x_{152}+x_{153}\rho,&g_1&=x_{156}+x_{157}\rho,
 &s_1&=x_{160}+x_{161}\rho,\\
 f_3&=C(x_{162}+x_{163}\rho),&g_3&=x_{166}+x_{167}\rho,
 &s_3&=x_{170}+x_{171}\rho,\\
 f_4&=C(x_{172}+x_{173}\rho),&g_4&=x_{176}+x_{177}\rho,
 &s_4&=x_{180}+x_{181}\rho,\\
 p&=x_{192}+x_{193}\rho.&&&&
\end{align*}
These symbols denote the displayed stored expressions; their
physical scalings are given by
\begin{equation}\label{eq:recovery}
 x(Q_i)=\frac{s_i^2}{q^2}+D\frac{f_i+g_i}{q}
 \quad(i=1,3,4),\qquad x(Q_2)=\frac p q.
\end{equation}
They are precisely the source point-recovery formulas after its
Laurent and fixed-parameter substitutions.

Define $A_0,A_1$ by the quadratic field expression
\begin{equation}\label{eq:components}
 A_0+A_1\rho
 =s_1^2+s_3^2+s_4^2
  +Dq(f_1+g_1+f_3+g_3+f_4+g_4)+qp.
\end{equation}
Then $A_j=q^2S_j$. In particular, both components have degree
at most two in the retained rational coordinates. No ordinate,
height or square coordinate is needed to recover $A_1$.

The physical detector arguments are
\[
 t_0=x_2/q,\qquad t_1=x_{86}/q,\qquad t_2=x_{200}/q.
\]
Use the two rows
\begin{align}
 2q(x_2-x_{86})+26q^2+A_0&=0,
   \label{eq:trace}\\
 2q(x_2-2x_{86}+x_{200})-8q^2+A_1&=0.
   \label{eq:shift}
\end{align}
The first replaces the source aggregate row, changing just its
$q^2$ coefficient from $14$ to $26$.
The second is the one additional aggregate row.
Dividing by the guarded $q^2$ shows that these rows are equivalent
to
\begin{equation}\label{eq:physical-shifts}
 t_1-t_0=13+S_0/2,\qquad
 t_2-t_1=17+(S_0-S_1)/2.
\end{equation}
The same conclusion follows by subtracting
\eqref{eq:shift} from~\eqref{eq:trace}; hence the two equations
control both components independently.

\section{The reflected contact equations}

Put
\[
 \Theta=\{0,-\tfrac12,-1,\tfrac1{2399},-\tfrac{32}{407}\},
 \qquad c=\tfrac3{11},\qquad h_c=\tfrac3{22},
 \qquad \Theta_c=\Theta\cup(c-\Theta).
\]
The reflection center $c$ belongs to the contact families discussed
in the pinned source. Its use here changes only the contact
equations; all scalar and pole inputs and all guards are retained.
Let
\[
 \alpha_\theta=(22\theta-3)^2,\qquad
 P_5(r)=\prod_{\theta\in\Theta}(r-\alpha_\theta)
       =r^5-e_1r^4+e_2r^3-e_3r^2+e_4r-e_5.
\]
The five paired constants are
\[
 9,\quad196,\quad625,\quad
 (7175/2399)^2,\quad(175/37)^2.
\]
Their elementary symmetric coefficients have common denominator
$L=7878870169$, with
\[
\begin{array}{c|r}
j&Le_j\\\hline
1&6786192246520\\
2&1229741066709366\\
3&42042541279847500\\
4&476802068222265625\\
5&1738195040039062500
\end{array}
\]

For each detector argument $t$ and its guarded inverse weight $w$,
the physical contact equations are
\begin{align*}
 r&=(22t-3)^2,&s&=r^2,\\
 v&=rs-e_1(1+w)s+e_2(1+w)r,&
 sv&=(1+w)(e_3s-e_4r+e_5).
\end{align*}
They are quadratic after the stored coordinate substitutions below.
Eliminating the defined $r,s,v$ gives
\[
 (1+w)P_5(r)=wr^5.
\]
Writing the physical weight as $q_t=1/w$ and using
\[
 P_5((22t-3)^2)
  =22^{10}\prod_{\theta\in\Theta}(t-\theta)(t-c+\theta)
\]
shows that this is exactly
\[
 (1+q_t)\prod_{\eta\in\Theta_c}(t-\eta)=(t-h_c)^{10}.
\]
The midpoint $t=h_c$ is included: it gives $q_t=w=-1$.
No division by $t-h_c$ is used. The contacts themselves cannot occur,
because none equals the midpoint and the right side is then nonzero.

For the final table put $q=x_{77}$. The stored arguments, squares,
fourth powers and combination coordinates have physical values
\[
\begin{array}{c|cccc|c}
j&t_j&r_j&s_j&v_j&w_j\\\hline
0&x_2/q&x_{74}/q&x_{75}/q&Cx_{76}/(Lq)&1/q\\
1&x_{86}/q&x_{87}/q&x_{88}/q&x_{89}/(Lq)&x_{91}/q\\
2&x_{200}/q&x_{201}/q&x_{202}/q&x_{203}/(Lq)&x_{205}/q
\end{array}
\]
The contact equations multiplied by $q^2$, or by $Lq^2$ as
appropriate, are precisely the twelve rows in the complete table.
Their degrees are at most two. The literal unit equations
$x_{77}x_{78}=1$, $x_{90}x_{91}=1$ and
$x_{204}x_{205}=1$ supply all required inverses.

The ten contacts are distinct and have diameter $25/11$.
The aggregate shifts are at least seven and nine. Consequently
the three translated ten-contact sets are pairwise disjoint.
Every square, inverse and realization coordinate used in these
equations is already included in the 229-coordinate count.

\section{Separation on the full point domain}

\begin{lemma}\label{lem:separation}
For every finite real point on $y^2=x^3-d^2x$ with $d>0$,
$x\ge-d$. Consequently the shifts in~\eqref{eq:physical-shifts}
satisfy
\[
 t_1-t_0\ge7,\qquad t_2-t_1\ge9.
\]
\end{lemma}
\begin{proof}
If $x<-d$, then all three factors of $x(x-d)(x+d)$ are negative,
contradicting its being a square. Apply this separately at the two
real embeddings of $\Q(\rho)$. The four-point sum therefore has
\[
 \sigma_+(\mathcal S)\ge-12+4\sqrt2,
 \qquad \sigma_-(\mathcal S)\ge-12-4\sqrt2.
\]
Their average gives $S_0\ge-12$. Moreover,
\[
 S_0-S_1=
 \frac{(2-\sqrt2)\sigma_+(\mathcal S)
 +(2+\sqrt2)\sigma_-(\mathcal S)}4\ge-16.
\]
Insert these two inequalities in~\eqref{eq:physical-shifts}.
\end{proof}

The contact set $\Theta_c$ has diameter $25/11$.
Lemma~\ref{lem:separation} therefore makes the three translates
\[
 \Theta_c,\quad \Theta_c-(t_1-t_0),\quad
 \Theta_c-(t_2-t_0)
\]
pairwise disjoint, including the finite torsion branches.
Each copied contact circuit has its own literal inverse guard,
and all guards are included in the coordinate and row counts.

The shifted sum $Z=\mathcal S+4d+1$ has positive value at both
real embeddings. The two aggregate equations give the exact identity
\[
 Z=-(2+2\rho)t_0+(2+4\rho)t_1-2\rho t_2-13+4\rho.
\]
All coefficients are algebraic integers, and the sum of their
absolute values is less than $32$ at either real embedding.
For the absolute logarithmic height this implies
\begin{equation}\label{eq:height}
 h(Z)\le h(t_0)+h(t_1)+h(t_2)+\log32.
\end{equation}
The statement follows place by place from the ultrametric inequality
at finite places and the triangle inequality at the two real places.

\section{Arithmetic inputs and the finite exceptional set}

The arithmetic proof uses the normalized elliptic group and point
interpretations in~\cite{ExplicitParent}, the prepared pole obstruction
and realization theorem stated there, its fixed-data simultaneous-prime
theorem and compatible index progressions, and its full-family
five-contact height theorem. We state the quantitative height input
precisely. Put
\[
 S=\{2,3,5,7,11,13,17,37,41,2399\},\qquad B_S=\prod_{p\in S}p,
\]
\[
 c_0=74712001/10^6,\qquad C_0=2^{68}3^{16}5^{16}7^4,
 \qquad E(R)=\prod_{p\mid R}(1+1/p).
\]
For $u=r/s$ in lowest terms, $s>0$, outside the four finite contacts
$0,1,2401,-343/32$, let $M_S(u)$ contain each prime outside $S$
having positive odd valuation on one of
\[
                       s,r,r-s,r-2401s,32r+343s.
\]
The retained full-family theorem, including its non-CM and CM branches,
is
\begin{equation}\label{eq:height-input}
 h(u)<c_0C_0 R^2E(R)\log R,\qquad R=B_SM_S(u).
\end{equation}
This is the theorem for this particular radical, identified as H4.1
in~\cite{ExplicitParent}. No reduction of an exceptional set is
inferred from a theorem with a different radical.

For the reflected contacts at center $c=3/11$, put
\[
 \Theta_c=\Theta\cup(c-\Theta),\qquad
 \Gamma=S\cup\{79,89,647,3593\}.
\]
Every contact denominator, every contact difference, and every
difference $c/2-\theta$ has prime support in $\Gamma$.
For $t\notin\Theta_c$, let $M(t)$ be the odd positive contact radical
outside $\Gamma$. The exact finite factorizations are supplied with
the coefficient and arithmetic certificates.

\subsection{Formal ratios without a numerical prime cutoff}

At an unramified odd place, with $r=v(z(Q))\ge1$ and $z=-x/y$,
the formal logarithm is
\[
 \ell(z)=\sum_{j\ge0}
 \frac{(-1)^j\binom{2j}{j}D^j}{4j+1}z^{4j+1}.
\]
When $p\ne5$, the first correction $-2Dz^5/5$ has valuation
at least $5r$. For $j\ge2$ the inequality
$v_p(4j+1)\le4j-4$ gives the same lower bound for every
subsequent term. The series converges and therefore
\begin{equation}\label{eq:local-log}
 v(\ell(Q))=v(z(Q)),\qquad
                  \ell(Q)/z(Q)\in1+p^{4r}\mathcal O.
\end{equation}
All primes outside $S$ satisfy these hypotheses and give good
reduction of the curve. This proof concerns the parameter $z$
itself; no higher-order logarithm truncation is required.

The complete source charts recover labels
\[
 n=m_1/2\in\tfrac12\Z,\quad a=m_2\in\Z\setminus\{0\},
 \quad t=m_3\in\Z,\quad b=2m_4+1\in2\Z+1.
\]
Put $c_a=|a|$ and $B=\max(1,|n|,|a|,|t|,|b|)$.
The physical scalar is
\[
 \beta=51X^2+2504,\quad
 r_4=\beta(\beta-1)(\beta-2401)(32\beta+343),\quad
 R_s=\frac{149^{2N+1}r_4}{1+149^2r_4}.
\]
The source's forced valuation at 149 excludes the reference torsion
and shifted-torsion strata before any logarithmic ratio is used.
In particular its reference index $h$ is nonzero and is not one,
and all reference ratios have their stated point interpretations.

\subsection{Scalar divisibility at all primes outside \texorpdfstring{$S$}{S}}

The four finite scalar contacts have collision support only
$\{2,3,5,7\}$. Outside $S$, a negative valuation of $\beta$
is $2v_p(X)$ and is even. Every contributing finite contact other
than 149 therefore gives an odd positive valuation
$\mu=v_p(R_s)$. At 149 the four residues
$(\theta-2504)/51$ are $59,97,109,108$, all nonsquares.
Thus 149 is absent from $M_S(\beta)$, while
$v_{149}(R_s)$ is independently $2N+1$ or $2N-1$.

The actual scalar equations give
\[
 W_s=Z_0^4/R_s^{10}+Z_1^2/R_s^5+F^2/R_s^4,
 \qquad b_s=R_s(R_s+9792)W_s^2,
\]
where $Z=2z(Q_2)/z(H)$ and $F=\operatorname{Nm}(1/x(H))\ne0$.
At every prime under consideration, $9792$ is a unit and $v_p(F)$
is even. The three finite term valuations in $R_s^4W_s$ have the
distinct classes $2+4\Z$, $1+2\Z$, and $4\Z$.
The pole obstruction forbids negative odd valuation of $b_s$.
Reading the three terms therefore gives
\[
 v_p(Z_0)>19\mu/8,\qquad v_p(Z_1)>9\mu/4,
 \qquad v_p(F)>7\mu/4.
\]
Zero components are omitted from this valuation comparison.

At some place above $p$, the reference is formal. The scalar point
is formal or reduces to $(0,0)$; translation by this two-torsion
point changes $z$ by a sign and makes it formal.
Multiplication by two kills the torsion terms in the index relation.
Equation~\eqref{eq:local-log} then gives
$v_p(a/h)>9\mu/4$. Since $h\in\tfrac12\Z$ is $p$-integral,
the same strict lower bound holds for $v_p(a)$.
Consequently, with
\[
 L_N=\left\lfloor\frac{9(2N-1)}4\right\rfloor+1,
\]
we obtain
\begin{equation}\label{eq:scalar-div}
              149^{L_N}M_S(\beta)^3\mid c_a.
\end{equation}
The H4.2 form of~\eqref{eq:height-input} has constant
$c_0C_0B_S^2E(B_S)\log B_S<2^{237}$.
Using $E(M)\le1+\log M$ and $(1+\log M)^2<2M$ gives
\begin{equation}\label{eq:scalar-height}
 \tau:=h(\beta)<2^{238}M_S(\beta)^3
                    \le2^{238}149^{-L_N}c_a.
\end{equation}

\subsection{The complete detector comparison}

At a positive odd contact of depth $\mu$ outside $\Gamma$, the
normalizing difference $t_j-c/2$ is a unit and the weight satisfies
$v_p(q_j)=-\mu$.
Write $C_i=L_i-L_i'$, $I=L_1-\beta$, and
\[
 G=L_3^2-(L_2+1)(L_2-1)L_3L_4+(L_2-1)^2L_4^2-1.
\]
The physical factors recovered by each complete detector are
\begin{align*}
 &1+q_j^5C_2^2+q_j^2D_H^2+q_j^4B_1^4,\\
 &1+q_j^5G^2+q_j^2A_3^2+q_j^4B_0^4,\\
 &1+q_j^5I^4+q_j^3A_2^2+q_j^2A_4^2+q_j^4B_3^4,\\
 &1+q_j^5C_3^2+q_j^2C_1^4+q_j^4B_4^4,\\
 &1+q_j^5C_4^2+q_j^2A_1^2+q_j^4B_2^4.
\end{align*}
Here $D_H=1/x(H)$, $A_i=1/x(U_i)$ and the $B_i$ are the
conjugate opposite reciprocals, including the reference predecessor
for $B_0$. Identity numerators have reciprocal zero.

All finite reciprocal valuations are even. The summands in each
factor have distinct valuation classes, so its valuation is at most
zero and no cancellation is possible. If $W_j$ is their product
and $\ell=\max_{v\mid p}v(W_j)\le0$, the trace equations give
\[
 b_j=\frac{(408/409)^2R_j}{\operatorname{Tr}(W_j^{-1})^2},
 \qquad v_p(b_j)\le\mu+2\ell.
\]
The formula is used here only at a consumed contact, where the
factor argument and the trace equation ensure that both denominators
are nonzero. Outside $\Gamma$, $v_p(408/409)\le0$; its negative
valuation at 409 only strengthens the inequality.
The pole obstruction gives $\ell>-\mu/2$. The individual terms
then imply
\[
 v(C_2,C_3,C_4,G)\ge3,\quad v(C_1,I)\ge1,
 \quad v(D_H,A_i,B_i,B_0)\ge2
\]
at a successful place.

One of $h$ and $h-1$ is a $p$-adic unit. The corresponding full
logarithmic ratio is integral; its positive comparison with the
other ratio transfers that integrality. Only after this step is
the relative precision in~\eqref{eq:local-log} used as an absolute
congruence. Multipliers 1, 2, 4 and the fourth-label offset 1
are integral, and 2 and 4 are units. This proves
\[
 L_1\equiv n\pmod p,\qquad
             (L_2,L_3,L_4)\equiv(a,t,b)\pmod{p^3}.
\]
The argument includes identity numerator points, two-torsion
differences, and split or inert places.

Define the integer residual
\[
 \Delta=t^2-(a+1)(a-1)tb+(a-1)^2b^2-1.
\]
If $n\ne\beta$, let $J_X$ be the absolute reduced numerator
of $n-\beta$. For each detector separately,
\begin{equation}\label{eq:residual-div}
 M(t_j)^3\mid\Delta\quad(\Delta\ne0),\qquad
 M(t_j)\mid J_X\quad(J_X\ne0).
\end{equation}
The elementary size bounds are
\[
                 |\Delta|\le14B^4,\qquad J_X\le4Be^\tau.
\]
No coprimality of the three radicals is asserted or needed.

\section{Effective soundness at \texorpdfstring{$N=15$}{N=15}}

Put
\[
 J_*=\prod_{p\in\Gamma}p
       =30366775457346956958330<2^{75},\qquad
 \mathcal C=c_0C_0J_*^2E(J_*)<2^{300}.
\]
For $u_+=2+1/t$ and $u_-=2+1/(3/11-t)$, the inverse
linear-fractional transformations have height constants at most
3 and 33. Outside $\Gamma$ their five-contact radicals are
disjoint and their product is $M=M(t)$. With
$R_\pm=B_SM_S(u_\pm)$, this gives
\[
 R_+R_-\le J_*^2M,\qquad
 E(R_+)E(R_-)\le E(J_*)^2E(M).
\]
Multiply~\eqref{eq:height-input} for the two arguments and take
square roots. The arithmetic-geometric mean inequality gives
\begin{equation}\label{eq:reflected-sharp}
 h(t)\le\log33+\mathcal C M\sqrt{E(M)}
                         (\log J_*+\tfrac12\log M).
\end{equation}
If $h(t)\le\log33$ the same conclusion is immediate.
This includes the midpoint $t=3/22$, where the height argument
does not divide by the normalizing difference.

The exact product of $1+1/p$ over primes $p\le1000$ outside
$\Gamma$ is less than $21/10$. For every larger prime,
$\log(1+1/p)<1/p<\log p/6006$. Hence
\begin{equation}\label{eq:euler}
                  E(M)<\frac{21}{10}M^{1/6006}.
\end{equation}
For a weaker estimate, use $E(R)\le1+\log R$ in the joint
height bound, $J_*<2^{75}$ and $c_0C_0<2^{149}$. It gives
\begin{equation}\label{eq:reflected-coarse}
       h(t)\le K M(22+\tfrac12\log M)^2,\qquad K=2^{304}.
\end{equation}
Indeed $1+75\log2<54<4\cdot22$; the factor 16 from squaring
and one further factor two absorb the additive $\log33$.

The four-point corollary in~\cite{TorsionSums}, followed by
\eqref{eq:height}, gives
\begin{equation}\label{eq:index-height}
                    B^2<40\sum_{j=0}^2h(t_j)+2425.
\end{equation}
It applies to every finite point recovered by the charts, with
arbitrary two-torsion translates and half-integral original indices.

Suppose $\Delta\ne0$. Apply
\eqref{eq:reflected-sharp}--\eqref{eq:euler} and the three separate
divisibilities~\eqref{eq:residual-div}. Absorbing the additive
constants gives
\[
 B^2<121\mathcal C\sqrt{21/10}\,
 |\Delta|^{1/3+1/36036}
                  (\log J_*+\tfrac16\log|\Delta|).
\]
Use $|\Delta|\le14B^4$ and write $s=\log B$.
Since $121\sqrt{21/10}\,14^{1/3+1/36036}<500$, we obtain
\begin{equation}\label{eq:delta-contradiction}
 e^{\frac23(1-1/6006)s}
                   <500\cdot2^{300}(53+2s/3).
\end{equation}
For $N\ge15$, scalar divisibility gives $s\ge66\log149$.
The quotient of the left side of~\eqref{eq:delta-contradiction}
by its final linear factor is increasing for $s\ge0$.
At the endpoint its logarithmic comparison is reversed:
\[
 \frac23(1-1/6006)66\log149
       >220\frac{6005}{6006}>219.963,
\]
whereas
\[
 \log500+300\log2+\log(53+44\log149)
 <\frac{25}{4}+300\frac{173287}{250000}+\frac{45}{8}
 =219.8194.
\]
Here $5<\log149<501/100$, $\log500<25/4$,
$\log2<173287/250000$ and $\log274<45/8$ have certified
rational series bounds. This contradiction proves $\Delta=0$.

The affine Lucas lemma in~\cite{ExplicitParent} now applies:
$b$ is odd and hence nonzero, and
\begin{equation}\label{eq:lucas-growth}
                   \log B\ge(c_a-6)\log(c_a-3).
\end{equation}
Its norm-one equation is obtained by putting $A=a+1$, $c'=|A|$
and $w=\operatorname{sgn}(A)(A-2)b$:
$t^2-c'tw+w^2=1$.
The integer-unit classification in $\Z[(c'+\sqrt{c'^2-4})/2]$
and the recurrence congruence modulo $c'\mp2$ force the Lucas
index to be at least $c'-2$. This gives~\eqref{eq:lucas-growth}
for $c_a\ge8$.

If $J_X\ne0$, use~\eqref{eq:reflected-coarse},
\eqref{eq:index-height} and $J_X\le4Be^\tau$. They imply
\[
 e^s<484K e^\tau(23+(s+\tau)/2)^2,
\]
and hence, by $\log(1+x)\le x$,
\[
 s<\frac{23}{22}\log(256036K)+\frac{12}{11}\tau<240+c_a.
\]
The last inequality uses~\eqref{eq:scalar-height} and
$149^{66}>2^{239}$, which give $\tau<c_a/2$.
But $c_a\ge149^{66}>600$ and $\log c_a>330$, so
\eqref{eq:lucas-growth} gives $s>320c_a>240+c_a$.
Thus $J_X=0$ and $\beta=n\in\tfrac12\Z$.
For $X=u/v$ in lowest terms, $v>0$, this gives $v^2\mid102$.
Squarefreeness of 102 proves $X\in\Z$.

Every numerical inequality strengthens as $N$ increases.
The exact arithmetic checker verifies the finite prime supports,
scalar nonsquare residues, product constants and decisive rational
logarithm bounds. The retained global arithmetic inputs and complete
witness construction remain mathematical proof obligations, proved
by their cited suppliers and the following argument.

\section{Complete witnesses for every integer input}

This construction requires only $N\ge3$; in particular it applies
at the sufficient soundness parameter below. Fix $X\in\Z$.
Then $r_4$ is a positive integer and
$R_s\ge149^{2N+1}/22202>2^{20}$.
The compact real neighborhoods proved in~\cite{ExplicitParent}
permit both $|z(H)|>1/2000$ and $|z(Q_2)|\le1$, with all reference
denominators nonzero. Hence $|Z_0|,|Z_1|<2^{12}$ and $|F|\le1/7$.
Uniformly in the input,
\[
              0<W_s<2^{49}R_s^{-4},\qquad
              0<b_s<2^{112}R_s^{-6}<17.
\]

Let $D_s=1+149^2r_4$. The numerators of $R_s$ and
$R_s+9792$ are coprime to $D_s$, so their product has square
reduced denominator. The scalar congruences give
$v_2(R_s)\in\{2,3\}$ and $v_3(R_s)=v_{17}(R_s)=0$;
$R_s(R_s+9792)$ is a square in $\Q_2$, $\Q_3$ and $\Q_{17}$.
First choose $H_*=[h_*]P_0$, with $64\mid h_*$, sufficiently
formal at the finite odd-valuation support, at 2, 3 and 17 whether
or not they occur in that support, and in the fixed real
neighborhoods. Next choose $a_0>0$, divisible by three, with
$Q_2=[a_0]P_0$ deeper at those places and close to the real identity.
The finite-index recurrence theorem permits $a_0$ to be arbitrarily
large. This gives square denominator and the required scalar
depths at 2, 3 and 17. After $a_0$ is fixed, continuity gives a
reference progression preserving these conditions.

Set $a=2a_0\in6\Z$. With $U_0=0$, $U_1=1$ and
$U_{k+1}=(a+1)U_k-U_{k-1}$, put
\[
       j=a-1,\qquad b=U_j/(a-1),\qquad t=U_{j+1}.
\]
The integral Lucas identities in~\cite{ExplicitParent} give
$b>1$ odd, $t>0$ and $\Delta=0$. Choose the four fixed points
\[
 Q_1=[\beta]P_0,\quad Q_2=[a_0]P_0,\quad
 Q_3=[t]R_0,\quad Q_4=[(b-1)/2]R_0.
\]
Their labels are exactly $\beta,a,t,b$. All are finite and
nontorsion. In particular both argument shifts in
\eqref{eq:physical-shifts} are now fixed before choosing contact
primes.

Let $U_0'=0$, $U_1'=t_1-t_0$, $U_2'=t_2-t_0$.
The thirty coefficients $\theta-U_j'$, for
$\theta\in\Theta_c$ and $0\le j\le2$, are distinct by the
shift separation. Choose an integer $D_0>0$ clearing their
denominators, including those of the midpoint shifts, by squares.
Before applying the prime theorem, shrink the reference neighborhoods
so that all tests are regular and bounded. Enlarge the finite
preparation set to include every fixed denominator, guard and shift
prime. At a prepared prime $p$, let $-B_p$ be a common lower
valuation bound, with $B_p\ge0$, for the finitely many powered test expressions in
the five factors. Prescribe the divisibility of $w$ sufficiently
deeply that $\lambda$ dominates all contact and midpoint shifts
and $v_p(q_j)\ge B_p+5$ for all three detectors. This makes every
nonconstant factor term have positive valuation, uniformly on those
neighborhoods. These data and the real bounds are fixed before the
analytic prime-selection parameter varies. Take the final reference
modulus $L$ to be a multiple of 64, the earlier scalar progression
modulus, and the orders of the relevant real component groups.
The fixed-data prime theorem gives
\[
 \lambda=u/(D_0^2w^2),\qquad
 p_{j,\theta}=u+D_0^2(U_j'-\theta)w^2
\]
simultaneously prime, outside every fixed exceptional and preparation
set, with the required Gaussian congruences and arbitrarily large
fixed prime and real-slope cutoffs. The selected primes may still be required to exceed 76176; soundness merely consumes
more primes than completeness needs to select.

The half-generator reduction and compatible progression theorem
then choose the final $H=[h]P_0$, with
$U_i=[h]Q_i$ and $V_i=[h-1]Q_i$, in the prescribed primary and
opposite formal groups while preserving all prepared and real
neighborhoods. All excluded chart denominators form a finite
exceptional set and are avoided. At a selected prime the labels
are integral and the absolute formal errors have order at least
four. Every nonconstant factor term therefore has positive
valuation at the successful place. The reciprocal trace is a
unit with residue one or two, giving valuation one for the actual
detector input $b_j$.

More explicitly, for each $j$ the contact ratio is
\[
 R_j=\frac{\prod_{\theta\in\Theta_c}p_{j,\theta}}
              {(u+D_0^2(U_j'-c/2)w^2)^{10}}.
\]
Its denominator is a square. Outside its ten selected numerator
primes every valuation is even, and division by the squared trace
preserves parity. Thus every actual $b_j$ has square reduced
denominator. At the fixed preparation places,
\[
 R_j=\prod_{\theta\in\Theta}
 \left(1-\frac{(c/2-\theta)^2}{(t_j-c/2)^2}\right)
\]
is arbitrarily close to one. The bounded test functions make
$W_j$ close to one and its trace close to two. The factor $408/409$
gives the required local squares and depths $4,2,2$ at 2, 3 and 17.
At both real embeddings, large positive slope gives $q_j>0$,
all factors positive, and
$b_j\longrightarrow(408/409)^2/4\in(0,17)$.

After this final reference is fixed, apply the prepared pole
realization theorem independently to $b_s,b_0,b_1,b_2$.
Their pole-coordinate sets are disjoint and include the complete
zero-auxiliary branches. Fill the reference and point charts by
their converses, and fill every product, square and contact helper
by its specified value. The three inverse weights are nonzero,
so all certified Laurent substitutions are defined. Both aggregate
rows hold by the chosen shifts. This constructs a rational zero
of the complete table for every integer input. No earlier pole
certificate is assumed to persist during the reference progression.

\section{Exact coefficient theorem}

Let $Q_{N,1},\ldots,Q_{N,207}$ be the specified integer quadratic
rows, and put
\[
 P_N=\sum_{j=1}^{207}Q_{N,j}^2.
\]

\begin{theorem}\label{thm:coefficients}
For every integer $N\ge1$, the polynomial $P_N$ has exactly 229
variables, degree four, content one, and 3,312 nonzero coordinate
monomials. Its exact coefficient height is
\[
 H(P_N)=\max\{9063965991649194484576286065455625302,149^{2N+2}\}.
\]
The constant term in this maximum dominates for $1\le N\le7$;
the parameter term dominates for $N\ge8$.
Over $\Q$, $P_N=0$ is equivalent to the complete three-detector
system with~\eqref{eq:physical-shifts}. Every square, height-recovery
expression and nonzero guard used in this system is included above.
\end{theorem}

\begin{proof}
The independent verifier authenticates the source table and checks
all 168 unchanged parent rows, the aggregate change, the six
changed parent contact rows, and the 34 coordinate copies. It checks
the 31 copied equations against the updated second detector and
independently reconstructs all twelve contact equations. Multiplication in
$\Q[\rho]/(\rho^2-2)$ independently verifies both components
of~\eqref{eq:components}. It reconstructs both aggregate rows and
their equivalence to~\eqref{eq:physical-shifts}. The three inverse
guards are literal rows numbered 26, 38 and 180.

Expansion by unordered pairs of terms gives 3,322 parameter-separated
terms and 3,312 coordinate-monomial profiles. Each profile is a
polynomial of degree at most two in $C$. After choosing its leading
sign, substitution $C=149+T$ gives positive constant coefficient and
nonnegative remaining coefficients. Thus none of these monomials
disappears for $C\ge149$. Every coordinate occurs.

Write $K=9063965991649194484576286065455625302$. A constant coefficient
profile of absolute value $K$ and the profile $22201C^2$ both occur. For $N=1,\ldots,7$,
exact evaluation gives height $K$. For every coefficient profile
$c_m(C)$, each of $22201C^2\pm c_m(C)$ has nonnegative coefficients
after substitution $C=149^8+T$. This proves the uniform height
formula. The coefficient of $X^4$ is one, proving both exact degree
four and content one. Finally a sum of rational squares vanishes
exactly when every summand vanishes.
\end{proof}

The coefficient crossover at $N=8$ concerns only this height formula.
The soundness and complete witness arguments prove the defining
threshold $N=15$, completing Theorem~\ref{thm:main}.

\begin{thebibliography}{9}
\bibitem{Parent195}
\emph{A quartic definition with 195 rational coordinates},
\nolinkurl{papers/quartic/quartic195_20261116_v45.tex},
with the complete table in
\nolinkurl{computation/quartic/affine_laurent195_v45/quadratic-family.json}.
\bibitem{TorsionSums}
\emph{Effective heights of elliptic sums with torsion translates},
17 November 2026,
\nolinkurl{papers/effective_sums/effective_torsion_sums_20261117_v62.tex}.
\bibitem{ExplicitParent}
\emph{$\Z$ is Diophantine in $\Q$}, Version 19, Part I,
17 November 2026. The three-detector construction there is attributed
to Gemini Carbon and builds on Jos\'e Brox's quadratic constructions.
Pinned commit
\texttt{498efba6fe75e9d28d6804f1391d7f70e4150afa},
\url{https://github.com/elliotglazer/diophantine-z-q}.
\end{thebibliography}
\end{document}
